\chapter{Contexto y estado del arte}
\section{Contexto del problema}
El fútbol es uno de los deportes con mayor impacto social y económico a nivel mundial. En los últimos años, los datos se han convertido en un campo más del deporte. Los sistemas de captura de eventos y de seguimiento óptico registran cada acción del partido con un nivel de detalle que hasta hace poco era inviable, y esa información se utiliza tanto para el análisis táctico como para la producción audiovisual o las apuestas deportivas \citep{goes2021bigdata, rein2016bigdata}.

En este contexto, la predicción del resultado de un partido constituye uno de los problemas más estudiados y más complejos de la analítica deportiva. El fútbol es un deporte de pocos goles, con una alta variabilidad entre encuentros y en el que el resultado final depende en gran medida de lo que ocurre durante el propio partido, de modo que las predicciones precisas resultan especialmente difíciles, incluso para profesionales del sector.

La mayor parte de los sistemas de predicción existentes se apoyan en información previa al inicio del encuentro, como resultados históricos, ratings de equipos o cuotas de apuestas. Sin embargo, estos enfoques no aprovechan la información que refleja la evolución del partido una vez iniciado, lo que limita su capacidad para adaptarse al estado real del juego. Tal y como se detalla en el estado del arte, algunos trabajos incorporan información temporal adicional, ya sea mediante actualizaciones dinámicas entre partidos o mediante el análisis secuencial de eventos intrapartido. No obstante, estos enfoques no siempre se orientan de forma explícita a la predicción del resultado final a medida que avanza el encuentro.

El objetivo del presente trabajo es abordar este problema desarrollando un sistema capaz de predecir la probabilidad de cada resultado (victoria local, empate y victoria visitante) a partir de datos históricos de eventos correspondientes a partidos ya disputados, transformados en métricas acumuladas minuto a minuto. De este modo, no se plantea un sistema de inferencia en tiempo real sobre un flujo de datos en vivo, sino una predicción dinámica intrapartido simulada a partir de información histórica, con aplicación potencial en contextos como el análisis táctico, el \textit{broadcasting} o las apuestas \textit{in-play}.

Como los modelos prepartido no capturan la evolución real del juego, este TFM propone usar eventos intrapartido agregados temporalmente para estimar cómo cambian las probabilidades de victoria local, empate y victoria visitante durante el encuentro.

\section{Estado del arte}
En los últimos años, la predicción de resultados en eventos deportivos ha sido abordada mediante diversas técnicas de aprendizaje automático, evolucionando desde modelos estadísticos tradicionales hacia enfoques más complejos basados en datos y aprendizaje profundo. En este contexto, los estudios existentes puede organizarse en diferentes líneas de investigación en función del tipo de datos utilizados y de las metodologías empleadas.

\subsection{Enfoques basados en datos históricos previos}

\subsubsection{Predictive analysis and modelling football results using machine learning approach for English Premier League}

El trabajo realizado por \cite{baboota2019predictive} aborda la predicción de resultados en partidos de fútbol de la Premier League (liga inglesa) mediante la aplicación de técnicas de \textit{machine learning}. El objetivo principal del estudio es analizar la capacidad de distintos algoritmos de clasificación para predecir el resultado final de un encuentro, formulado como un problema multiclase con tres posibles salidas: victoria local, empate o victoria visitante.

Para ello, los autores utilizan datos históricos de partidos, a partir de los cuales construyen un conjunto de variables representativas del rendimiento de los equipos. Estas variables incluyen estadísticas agregadas como resultados previos, rachas de victorias o derrotas y métricas derivadas del desempeño reciente.

El estudio evalúa diferentes algoritmos de aprendizaje supervisado, entre los que destacan la regresión logística, los árboles de decisión y los métodos de ensamblado, en particular Random Forest. Los resultados obtenidos muestran que los modelos basados en ensamblado superan a los modelos lineales tradicionales en términos de precisión, alcanzando valores en torno al 65--70\%. Este resultado pone de manifiesto la capacidad de los modelos no lineales para capturar relaciones complejas entre las variables de entrada y el resultado del partido.

No obstante, el enfoque presenta varias limitaciones. En primer lugar, el modelo se basa exclusivamente en información previa al partido, lo que impide capturar la evolución dinámica del juego y los eventos que ocurren durante el encuentro. En segundo lugar, las variables utilizadas son principalmente agregadas, lo que reduce la capacidad del modelo para representar la naturaleza secuencial y contextual del fútbol. Por último, el estudio no aborda en profundidad la interpretabilidad del modelo, limitando la comprensión del impacto de cada variable en la predicción final.

A pesar de estas limitaciones, este trabajo constituye una base sólida en la aplicación de técnicas de aprendizaje automático a la predicción de resultados deportivos, demostrando la viabilidad de estos enfoques y estableciendo un punto de partida para investigaciones posteriores orientadas a la incorporación de datos en tiempo real y modelos más complejos.

\subsubsection{Random forest model identifies serve strength as a key predictor of tennis match outcome}

El trabajo desarrollado por \cite{gao2019random} aborda la predicción de resultados en partidos de tenis mediante el uso de técnicas de \textit{machine learning}, con especial enfoque en modelos basados en \textit{Random Forest}. El objetivo principal del estudio es evaluar la capacidad de estos modelos para predecir el resultado de un partido a partir de datos históricos detallados de encuentros y jugadores.

Para ello, los autores recopilan y procesan una de las bases de datos más completas de partidos de tenis disponibles, realizando tareas de limpieza y preparación de datos con el fin de construir un conjunto de variables representativas del rendimiento de los jugadores. Estas variables incluyen métricas relacionadas con el desempeño en el servicio, puntos ganados y otras estadísticas relevantes del juego.

El enfoque metodológico se basa en la aplicación de modelos relativamente simples de aprendizaje automático, lo que permite tiempos de entrenamiento y predicción reducidos, facilitando la incorporación de nuevos datos y la posibilidad de realizar predicciones de manera eficiente. A partir de estos modelos, el estudio alcanza niveles de precisión superiores al 80\%, superando incluso las predicciones basadas en cuotas de apuestas tradicionales.

Uno de los principales aportes del trabajo es la identificación de la fuerza del servicio como la variable más determinante en la predicción del resultado del partido. En particular, métricas como el porcentaje de puntos ganados con el primer y segundo servicio se revelan como factores clave en el rendimiento de los jugadores y, por tanto, en el desenlace del partido. Este descubrimiento pone de manifiesto la importancia de analizar la contribución individual de las variables en modelos predictivos deportivos.

El estudio presenta limitaciones relevantes. Al igual que otros trabajos basados en datos históricos, el modelo no incorpora información en tiempo real ni considera la evolución del partido durante su desarrollo, lo que limita su aplicabilidad en escenarios dinámicos. Asimismo, aunque identifica variables clave, el análisis de interpretabilidad se limita a la importancia de variables proporcionada por el modelo, sin explorar técnicas más avanzadas de explicación.

A pesar de estas limitaciones, este trabajo destaca por demostrar que modelos relativamente simples pueden alcanzar altos niveles de precisión en la predicción de resultados deportivos, así como por su contribución en la identificación de variables con alto poder predictivo, aspecto especialmente relevante en el contexto del presente trabajo.

\subsubsection{The use of machine learning in sport outcome prediction: A review}

El trabajo desarrollado por \cite{horvat2020review} presenta una revisión exhaustiva del uso de técnicas de \textit{machine learning} en la predicción de resultados deportivos. En este estudio, los autores analizan más de 100 trabajos previos con el objetivo de identificar las metodologías más empleadas, los tipos de datos utilizados y las principales tendencias en este campo de investigación.

Uno de los principales hallazgos del estudio es que la mayoría de los trabajos abordan la predicción de resultados deportivos como un problema de clasificación supervisada, donde el objetivo es predecir una única salida (por ejemplo, victoria o derrota), siendo menos frecuente el enfoque basado en predicción de valores numéricos. Asimismo, se observa que los algoritmos más utilizados incluyen redes neuronales, máquinas de soporte vectorial y métodos basados en árboles de decisión, destacando el uso recurrente de modelos no lineales para capturar la complejidad de los datos deportivos.

En relación con los datos, el estudio pone demuestra que prácticamente todos los trabajos analizados emplean algún tipo de selección y extracción de características antes de aplicar los modelos de aprendizaje automático. Este proceso resulta fundamental para mejorar el rendimiento predictivo, dado que permite identificar las variables más relevantes dentro de un contexto caracterizado por alta variabilidad y ruido. Además, los autores señalan que los conjuntos de datos suelen organizarse de forma cronológica, utilizando técnicas como la segmentación temporal y la validación cruzada para evaluar el rendimiento de los modelos.

Otro aspecto relevante identificado en la revisión es la falta de estandarización en los enfoques utilizados. La diversidad en los conjuntos de datos, las variables empleadas y las metodologías de evaluación dificulta la comparación directa entre estudios y limita la generalización de los resultados obtenidos. Asimismo, los autores destacan que, a pesar del uso de técnicas avanzadas de aprendizaje automático, la precisión de los modelos se ve condicionada por la naturaleza inherentemente incierta de los eventos deportivos.

Finalmente, el estudio concluye que, aunque el uso de \textit{machine learning} ha permitido avances significativos en la predicción de resultados deportivos, todavía existen importantes retos en el campo. Entre ellos, destacan la necesidad de mejorar la calidad y consistencia de los datos, así como la incorporación de nuevas fuentes de información que permitan capturar mejor la complejidad de los eventos deportivos.

Este trabajo proporciona una visión global del estado actual de la investigación en predicción deportiva mediante \textit{machine learning}, permitiendo contextualizar los estudios previamente analizados y evidenciando tanto las fortalezas como las limitaciones de los enfoques basados en datos históricos.

\subsubsection{Síntesis de los estudios revisados}

Para facilitar la comparación entre los trabajos analizados, se recoge en la Tabla \ref{tab:sintesis_historicos} una síntesis de los principales elementos de cada estudio. Esta tabla permite observar de forma conjunta el tipo de datos utilizado, los modelos aplicados, los resultados reportados y la principal relación de cada trabajo con el presente TFM.

\begin{table}[H]
    \centering
    \scriptsize
    \renewcommand{\arraystretch}{1.25}
    \begin{tabularx}{\textwidth}{|p{2.5cm}|p{2.7cm}|p{2.7cm}|p{2.5cm}|X|}
        \hline
        \textbf{Estudio}                                                                                                                         & \textbf{Dataset} & \textbf{Modelo o enfoque} & \textbf{Métricas o resultados} & \textbf{Aportación y relación con este TFM} \\
        \hline

        Baboota y Kaur (2019)                                                                                                                    &
        Partidos históricos de la Premier League, con variables agregadas del rendimiento previo de los equipos.                                 &
        Regresión logística, árboles de decisión y Random Forest.                                                                                &
        Precisión aproximada del 65--70\%, con mejor rendimiento de los modelos de ensamblado.                                                   &
        Referencia directa para formular la predicción de resultados de fútbol como un problema de clasificación multiclase. Su principal limitación es que trabaja con información previa al partido y no incorpora eventos intrapartido.                                     \\
        \hline

        Gao y Kowalczyk (2019)                                                                                                                   &
        Datos históricos de partidos de tenis, con estadísticas detalladas de jugadores y partidos.                                              &
        Random Forest.                                                                                                                           &
        Precisión superior al 80\% en la predicción del resultado del partido.                                                                   &
        Aunque se aplica a otro deporte, muestra que los modelos de ensamblado pueden obtener buenos resultados cuando las variables representan correctamente el rendimiento deportivo. También refuerza la importancia de analizar la contribución de las variables.         \\
        \hline

        Horvat y Job (2020)                                                                                                                      &
        Revisión de más de 100 estudios sobre predicción de resultados deportivos.                                                               &
        Análisis comparativo de técnicas de machine learning, como redes neuronales, máquinas de soporte vectorial y modelos basados en árboles. &
        No presenta una métrica única, al tratarse de una revisión de distintos estudios.                                                        &
        Contextualiza la predicción deportiva como un problema habitual de clasificación supervisada y destaca la importancia de la selección de características, la calidad de los datos y la falta de estandarización entre estudios.                                        \\
        \hline
    \end{tabularx}
    \caption{Síntesis de trabajos basados en datos históricos previos.}
    \label{tab:sintesis_historicos}
\end{table}

\subsubsection{Comparación entre estudios}

Los estudios revisados muestran distintas formas de abordar la predicción deportiva a partir de datos históricos. El trabajo de Baboota y Kaur (2019) es el más directamente relacionado con este TFM, ya que plantea la predicción de resultados de fútbol como un problema de clasificación multiclase con tres posibles salidas: victoria local, empate y victoria visitante. Además, compara modelos clásicos como regresión logística, árboles de decisión y Random Forest, observando que los modelos de ensamblado ofrecen mejor rendimiento que los enfoques más simples.

Por su parte, el trabajo de Gao y Kowalczyk (2019), aunque pertenece al ámbito del tenis, resulta útil porque refuerza una idea también aplicable al fútbol: el rendimiento del modelo depende en gran medida de la calidad de las variables utilizadas. En este caso, Random Forest obtiene buenos resultados porque las variables de entrada describen aspectos muy relevantes del juego, especialmente el rendimiento en el servicio. En comparación con Baboota y Kaur, este estudio no se centra en fútbol ni en un problema de tres clases, pero sí confirma la utilidad de los modelos de ensamblado y del análisis de importancia de variables en predicción deportiva.

La revisión de Horvat y Job (2020) aporta una visión más general. A diferencia de los dos estudios anteriores, no propone un modelo concreto ni un único conjunto de datos, sino que analiza tendencias comunes en la literatura sobre predicción deportiva. Esta revisión permite contextualizar los resultados de Baboota y Kaur y de Gao y Kowalczyk, mostrando que la clasificación supervisada, la selección de características y los modelos no lineales son enfoques habituales en este tipo de problemas.

En conjunto, estos trabajos muestran que los modelos clásicos de machine learning pueden ser una base sólida para la predicción deportiva, especialmente cuando se dispone de variables bien construidas. Sin embargo, también presentan una limitación común: se basan principalmente en información histórica o previa al partido, por lo que no capturan directamente cómo cambia el estado del encuentro durante su desarrollo. Esta limitación justifica que el presente TFM incorpore eventos intrapartido agregados temporalmente como entrada del modelo predictivo.

\subsection{Enfoques basados en información temporal y dinámica del juego}

\subsubsection{Using ELO ratings for match result prediction in association football}

El trabajo desarrollado por \cite{hvattum2010elo} analiza el uso del sistema de puntuación ELO para la predicción de resultados en fútbol. El objetivo principal del estudio es evaluar si las valoraciones dinámicas de los equipos, actualizadas tras cada partido, pueden mejorar la precisión en la estimación de resultados frente a otros enfoques tradicionales.

El sistema ELO, originalmente desarrollado para el ajedrez, asigna a cada equipo una puntuación que refleja su rendimiento relativo. Esta puntuación se actualiza tras cada encuentro en función del resultado obtenido y de la diferencia esperada entre los equipos. De este modo, el modelo incorpora de manera implícita la evolución temporal del rendimiento, adaptándose a cambios en la forma de los equipos a lo largo de la competición.

En el estudio, los autores utilizan las puntuaciones ELO como variable principal dentro de modelos de regresión para predecir el resultado de los partidos. Los resultados obtenidos muestran que el uso de ratings dinámicos mejora la capacidad predictiva en comparación con modelos basados únicamente en estadísticas históricas agregadas. Además, el enfoque presenta la ventaja de ser computacionalmente eficiente y fácilmente actualizable tras cada partido, lo que lo hace adecuado para entornos donde se requiere una actualización continua de las predicciones.

No obstante, el modelo presenta limitaciones relevantes. En primer lugar, aunque incorpora una dimensión temporal, esta se basa únicamente en resultados finales de partidos, sin considerar eventos internos del juego como tiros, posesión o acciones específicas. En segundo lugar, el sistema ELO simplifica la complejidad del fútbol al representar el rendimiento de un equipo mediante un único valor escalar, lo que limita su capacidad para capturar dinámicas más complejas del juego.

\subsubsection{Actions Speak Louder than Goals: Valuing Player Actions in Soccer}

El estudio realizado por \cite{decroos2019actions} propone un enfoque avanzado para el análisis del rendimiento en fútbol basado en la valoración de acciones individuales durante el partido. A diferencia de los modelos tradicionales centrados en estadísticas agregadas o resultados finales, este estudio utiliza secuencias de eventos del juego, como pases, tiros, recuperaciones y pérdidas de balón, para modelar el desarrollo del encuentro de manera más precisa.

El núcleo del trabajo es el modelo denominado VAEP (\textit{Valuing Actions by Estimating Probabilities}), que asigna un valor a cada acción en función de cómo modifica la probabilidad de que un equipo marque o conceda un gol en un futuro cercano. Para ello, se emplean modelos de aprendizaje automático que estiman estas probabilidades a partir del estado actual del juego, teniendo en cuenta el contexto espacial y temporal de cada acción.

Este enfoque permite analizar el partido como una secuencia dinámica de eventos, capturando la evolución del juego de forma secuencial y proporcionando una representación más detallada del impacto de cada acción. De este modo, el modelo no se limita a evaluar resultados finales, sino que ofrece una valoración continua del estado del partido a medida que se producen los eventos.

Uno de los principales aportes del estudio es la capacidad de cuantificar la contribución individual de cada acción al resultado del partido, lo que permite identificar patrones de juego y evaluar el rendimiento de jugadores y equipos de forma más granular. Además, el uso de datos secuenciales y contexto espacial supone un avance significativo respecto a enfoques basados únicamente en estadísticas agregadas.

Sin embargo, el modelo presenta limitaciones. En particular, su complejidad computacional es superior a la de modelos más simples, y requiere acceso a datos detallados de eventos, lo que puede dificultar su implementación en algunos contextos. Asimismo, aunque proporciona una estimación continua del impacto de las acciones, no está diseñado directamente para predecir el resultado final del partido en términos de victoria, empate o derrota.

Pese a estas limitaciones, este trabajo representa un avance fundamental en el análisis del fútbol mediante \textit{machine learning}, al introducir un enfoque basado en eventos que permite modelar la evolución del juego de manera dinámica.

\subsubsection{Sports match prediction model based on LSTM and attention mechanism}

La investigación desarrollada por \cite{zhang2022lstm} propone un modelo de predicción de resultados deportivos basado en redes neuronales recurrentes, concretamente en arquitecturas Long Short-Term Memory (LSTM) combinadas con mecanismos de atención. El objetivo principal del estudio es mejorar la precisión en la predicción de resultados mediante la modelización de la información como una serie temporal.

A diferencia de los enfoques tradicionales basados en estadísticas agregadas, este modelo considera la evolución temporal del rendimiento de los equipos, introduciendo los datos de partidos en una secuencia ordenada en el tiempo. De este modo, el modelo es capaz de capturar dependencias temporales y patrones dinámicos en el comportamiento de los equipos, lo que resulta fundamental en un contexto tan variable como el deportivo.

El modelo propuesto, denominado AS-LSTM, incorpora un mecanismo de atención que permite asignar mayor peso a ciertos momentos o características dentro de la secuencia temporal, mejorando así la capacidad del modelo para identificar información relevante. Este enfoque permite no solo capturar la evolución del rendimiento, sino también diferenciar qué eventos o periodos tienen mayor impacto en la predicción final.

Los resultados obtenidos muestran que el modelo basado en LSTM con atención supera a modelos tradicionales de aprendizaje automático, demostrando una mayor capacidad para capturar patrones complejos en datos secuenciales y mejorar la precisión en la predicción de resultados deportivos.

Aun así, el enfoque presenta algunas limitaciones. En primer lugar, requiere grandes volúmenes de datos para su entrenamiento, lo que puede dificultar su aplicación en contextos con disponibilidad limitada de información. Además, la complejidad del modelo reduce su interpretabilidad en comparación con modelos más simples, dificultando la comprensión de las decisiones tomadas por el sistema.

\subsubsection{Síntesis y enfoque del TFM}

Los trabajos revisados en este apartado muestran distintas formas de incorporar la dimensión temporal en la predicción deportiva. El sistema ELO introduce una actualización dinámica del rendimiento de los equipos, pero lo hace a partir de resultados anteriores y no mediante los eventos internos del partido. Por otro lado, enfoques como VAEP utilizan datos de eventos para valorar la contribución de cada acción, aunque su objetivo principal no es predecir directamente el resultado final del encuentro en términos de victoria local, empate o victoria visitante. Finalmente, los modelos secuenciales basados en LSTM y mecanismos de atención permiten representar dependencias temporales más complejas, pero suelen requerir grandes volúmenes de datos y presentan una menor interpretabilidad.

A partir de estas limitaciones, el presente TFM adopta un enfoque basado en eventos intrapartido agregados temporalmente. En lugar de modelar la secuencia completa de acciones mediante redes neuronales recurrentes, se transforman los eventos registrados durante el partido en un conjunto de métricas temporales que representan el estado del encuentro en distintos momentos del juego.

Este planteamiento busca aprovechar la información dinámica que aportan los eventos sin asumir la complejidad de un modelo secuencial profundo. Por ello, el trabajo se plantea como una predicción dinámica intrapartido simulada a partir de datos históricos, cuyo objetivo es estimar la probabilidad de victoria local, empate y victoria visitante utilizando únicamente la información disponible hasta cada instante analizado.

\section{Conclusiones del estado del arte}

La revisión realizada permite extraer varias conclusiones relevantes para orientar el desarrollo del presente trabajo. En primer lugar, los enfoques basados en datos históricos previos muestran que las técnicas de machine learning pueden ser útiles para la predicción de resultados deportivos, especialmente cuando se dispone de variables representativas del rendimiento de los equipos. Trabajos como el de Baboota y Kaur (2019) demuestran que la predicción del resultado de un partido de fútbol puede formularse como un problema de clasificación multiclase, mientras que otros estudios refuerzan la importancia de la selección de características y del uso de modelos capaces de capturar relaciones no lineales.

Sin embargo, estos enfoques presentan una limitación clara: suelen trabajar con información disponible antes del inicio del partido o con estadísticas agregadas de encuentros anteriores. Esto hace que no puedan adaptarse al desarrollo real del juego una vez el partido ha empezado. En un deporte como el fútbol, donde el marcador, el ritmo del partido y las acciones relevantes pueden cambiar de forma importante la probabilidad de cada resultado, esta limitación reduce la utilidad de los modelos puramente prepartido.

Por otro lado, los enfoques basados en información temporal y dinámica del juego muestran que incorporar la evolución del partido puede aportar una representación más rica del problema. Modelos como ELO introducen una actualización dinámica del rendimiento, mientras que enfoques basados en eventos, como VAEP, permiten valorar la contribución de las acciones individuales. A su vez, los modelos secuenciales como las redes LSTM permiten capturar dependencias temporales más complejas. No obstante, estos enfoques también presentan limitaciones: algunos no consideran los eventos internos del partido, otros no predicen directamente el resultado final en términos de victoria, empate o derrota, y los modelos secuenciales profundos pueden requerir una mayor complejidad de diseño, entrenamiento e interpretación.

A partir de estas conclusiones, este trabajo se sitúa en un punto intermedio entre los modelos basados únicamente en información histórica previa y los modelos secuenciales más complejos. El enfoque adoptado se basa en utilizar eventos intrapartido de partidos históricos y transformarlos en métricas temporales que representen el estado del encuentro en distintos momentos del juego. De esta forma, se pretende aprovechar la información dinámica del partido sin asumir la complejidad de modelar directamente toda la secuencia de eventos.

Por tanto, este TFM propone desarrollar un sistema que permita predecir cómo cambian las probabilidades de victoria local, empate y victoria visitante a medida que avanza un partido. Para ello, se utilizarán datos históricos de eventos y se tendrá en cuenta únicamente la información disponible hasta cada momento analizado. Con esta propuesta se busca cubrir una limitación detectada en los trabajos revisados: muchos modelos predicen el resultado antes del partido, pero no aprovechan de forma directa lo que ocurre durante el juego.